\clearpage
\fancyhead[L]{\small CERES ELEMENTE XII.}
\hypertarget{native-304-lower}{}
\begin{center}CERES ELEMENTE XII.
\end{center}

\begin{center}
\fontsize{10}{12}\selectfont\begin{tabular}{c}
\hline
Göttingische gelehrte Anzeigen. Stück 40. Seite 393..395. 1808 März 10. \\
\hline
\end{tabular}
\end{center}

Die Ausbeute der vorigjährigen Beobachtungen der neuen Planeten, Vesta, Pallas und Juno, ist
bereits in den Händen des astronomischen Publicums: die die Juno betreffenden Resultate wurden
erst neuerlich in diesen Blättern mitgetheilt. Es ist also nur noch die Ceres übrig, über welche
Hr. Prof. \textsc{Gauss} seine Untersuchungen so eben geendigt hat. Ausser den vom Hrn. Prof.
\textsc{Harding} im April und Mai 1807 auf der hiesigen Sternwarte angestellten und einigen
Berliner Beobachtungen konnte Hr. Prof. \textsc{Gauss} noch eine schöne Reihe handschriftlich
mitgetheilter Mailänder Beobachtungen benutzen, aus welchen allen sich für die letzte Opposition
folgendes Resultat ergab:


\begin{center}\fontsize{10}{12}\selectfont\setlength{\tabcolsep}{3pt}\renewcommand{\arraystretch}{1.15}\begin{tabular}{lr}
Zeit der Opposition 1807, Mai 3. .\;.\; & $4^{\text{u}}\,12^{\text{m}}\,44^{\text{s}}$\\
\multicolumn{2}{r}{Meridian von Göttingen}\\
Wahre Länge .\;.\;.\;.\;.\;.\;.\;,\;.\; & $222^\circ\;14'\;9''1$\\
Wahre geocentrische Breite .\;.\;.\;.\; & $10\;40\;15.9$ nordlich.\\
\end{tabular}\end{center}

\noindent Dies ist die \emph{fünfte} Opposition, die seit Entdeckung der \emph{Ceres} beobachtet ist.

\par\indent Die letzten, auch in diesen Blättern (1806 S. 1946) [von \textsc{Harding}] mitgetheilten,
Elemente weichen zwar von diesem Resultate nur 3 Minuten in der Länge, und nur 13 Secunden in der
heliocentrischen Breite ab; inzwischen gibt jener Unterschied, so klein er ist, doch einen
entscheidenden Beweis, dass diejenigen Störungen, worauf Hr. Prof. \textsc{Gauss} sich bisher
eingeschränkt hat, nicht mehr zureichen, alle vorhandenen Beobachtungen genau darzustellen:
bisher sind nemlich blos die vom Jupiter herrührenden Störungen in Betrachtung gezogen, und auch
hier diejenigen ausgeschlossen, welche von den höheren Potenzen der Excentricitäten abhängen. Um
also die bisherigen und künftigen Beobachtungen genau vereinigen zu können, werden besonders die
letzten Störungen nachzuholen sein, obwohl es aus mehr als einer Ursache rathsam sein wird, zu
dieser Arbeit noch eine oder ein paar neue Oppositionen abzuwarten. Diesmal hat also Hr. Prof.
\textsc{Gauss} sich noch begnügt, blos mit Zuziehung derjenigen Störungen der Länge und des
Radius Vectors, die im Märzhefte der Monatlichen Correspondenz von 1803 in Tafeln gebracht sind,
jedoch mit vollständiger berechneten Breitenstörungen (welche bisher noch nicht gedruckt worden),
die letzten Elemente so zu verbessern, dass die Differenzen unter alle fünf Oppositionen nach
Möglichkeit vertheilt werden. Das Resultat dieser Untersuchung sind folgende neue Elemente (der
Zahl nach Nr. XII), die, wie man bemerken wird, von den letzten nur sehr wenig verschieden sind.
Die Epochen gelten für den Meridian von Göttingen.

\Needspace{18\baselineskip}
\SourceLeaf{305}
\begin{center}\fontsize{10}{12}\selectfont\setlength{\tabcolsep}{3pt}\renewcommand{\arraystretch}{1.15}\begin{tabular}{lrr}
& Mittlere Länge & Sonnennähe\\\hline
1802 & $155^\circ\;28'\;15''7$ & $146^\circ\;30'\;22''1$\\
1803 & $233\;38\;7.1$ & $32\;23.3$\\
1804 & $312\;0\;49.5$ & $34\;24.9$\\
1805 & $30\;10\;40.9$ & $36\;26.1$\\
1806 & $108\;20\;32.3$ & $38\;27.4$\\
1807 & $186\;30\;23.7$ & $40\;28.6$\\
1808 & $264\;53\;6.1$ & $42\;30.2$\\
\end{tabular}\end{center}

\begin{center}\fontsize{10}{12}\selectfont\setlength{\tabcolsep}{3pt}\renewcommand{\arraystretch}{1.15}\begin{tabular}{lr}
Tägliche mittlere Bewegung .\;.\;.\;.\;.\; & $770''9354$\\
Excentricität 1806 .\;.\;.\;.\;.\;.\;.\;.\;.\; & $0.0785251$\\
Jährliche Abnahme .\;.\;.\;.\;.\;.\;.\;.\; & $0.00000583$\\
Logarithm. d. halben grossen Axe .\;.\;.\; & $0.4420439$\\
Aufsteigender Knoten 1806 .\;.\;.\;.\; & $80^\circ\;53'\;24''6$\\
Jährliche Bewegung .\;.\;.\;.\;.\;.\;.\;.\; & $+1''48$\\
Neigung der Bahn 1806 .\;.\;.\;.\;.\;.\; & $10^\circ\;37'\;33''2$\\
Jährliche Abnahme .\;.\;.\;.\;.\;.\;.\;.\; & $-0''44$\\
\end{tabular}\end{center}

\par\indent Mit den fünf bisherigen Oppositionen stimmen diese Elemente folgender Massen überein:

\begin{center}\fontsize{9}{11}\selectfont\setlength{\tabcolsep}{3pt}\renewcommand{\arraystretch}{1.15}\begin{tabular}{lrrrr}
& \multicolumn{2}{c}{Unterschied} & \multicolumn{2}{c}{\shortstack{Nach den\\XIII. Elementen}}\\
& \shortstack{heliocentr. Länge\\in der Bahn} & \shortstack{heliocentrische\\Breite} & & \\\hline
1802 & $-24''1$ & $+1''2$ & $-0''5$ & $-1''9$\\
1803 & $+40.9$ & $-9.6$ & $+22.1$ & $-5.3$\\
1804 & $-21.0$ & $-3.1$ & $-34.9$ & $-1.2$\\
1806 & $+7.8$ & $0$ & $+26.7$ & $-3.6$\\
1807 & $-9.5$ & $-11.3$ & $-26.9$ & $-11.7$\\
1808 & $+62.7$ & & $+20.8$ & $-6.0$\\
\end{tabular}\end{center}

\noindent [Die Vergleichung mit der Opposition v. 1808 und die Unterschiede für die Elemente Nr. XIII. ist
von \textsc{Gauss} handschriftlich beigefügt. \textsc{Schering}.]

\Needspace{20\baselineskip}
\fancyhead[L]{\small COMET VON 1807.}
\begin{center}
COMET VON 1807.
\end{center}

\begin{center}
\fontsize{10}{12}\selectfont\begin{tabular}{c}
\hline
Göttingische gelehrte Anzeigen. Stück 53. Seite 521..523. 1808 April 2. \\
\hline
\end{tabular}
\end{center}

Hr. Inspector \textsc{Bessel} in Lilienthal ist, einem vom 18. März datirten Schreiben an Hrn.
Professor \textsc{Gauss} zufolge, so glücklich gewesen, den letzten Cometen noch bis gegen Ende
des Februars zu verfolgen. Wir eilen, hier diese schätzbaren Beobachtungen mitzutheilen.


\begin{center}\fontsize{10}{12}\selectfont\setlength{\tabcolsep}{3pt}\renewcommand{\arraystretch}{1.15}\begin{tabular}{lrrr}
1808 & Mittlere Zeit & \shortstack{Scheinbare\\ger. Aufsteigung} & \shortstack{Scheinbare\\nordl. Abweich.}\\\hline
Jan. 4 & $9^{\text{u}}\,49^{\text{m}}\,40^{\text{s}}$ & $325^\circ\;17'\;16''1$ & $47^\circ\;13'\;11''1$\\
4 & $11\;34\;32$ & $325\;21\;47.6$ & $47\;13\;40.9$\\
12 & $6\;58\;59$ & $333\;3\;58.2$ & $47\;40\;58.8$\\
21 & $7\;26\;15$ & $341\;19\;13.0$ & $47\;59\;15.3$\\
23 & $6\;49\;6$ & $343\;2\;37.9$ & $48\;2\;5.2$\\
Febr. 19 & $7\;20\;6$ & $3\;43\;35.7$ & $48\;21\;17.1$\\
20 & $7\;23\;15$ & $4\;24\;3.7$ & $48\;21\;33.5$\\
24 & $7\;35\;16$ & $7\;3\;11.1$ & $48\;23\;38.3$\\
\end{tabular}\end{center}

\par\indent Unter den im Februar beobachteten Declinationen ist, nach Hrn. \textsc{Bessel}'s eigenem Urtheile,
die vom 20. die einzige, die Vertrauen verdient.

W\SourceLeaf{306}ir fügen diesen Beobachtungen auch noch die späteste, von Hrn. Dr. \textsc{Olbers} in Bremen
gemachte, bei; die frühern dieses Astronomen sind bereits durch die Monatl. Correspondenz
bekannt geworden:

\begin{quote}
Febr. 14. $7^{\text{u}}\,30^{\text{m}}\,36^{\text{s}}$ mittl. Z.\quad ger. Aufsteig. $0^\circ\;15'\;53''$,\quad nordl. Abweichung $48^\circ\;18'\;10''$
\end{quote}

\par\indent Hr. \textsc{Bessel} hat diese spätesten Beobachtungen benutzt, um seine parabolischen Elemente
nochmals zu verbessern, und folgende Resultate gefunden, die von den Elementen des Herrn
Prof. \textsc{Gauss} wenig verschieden sind:


\begin{center}\fontsize{9}{11}\selectfont\setlength{\tabcolsep}{3pt}\renewcommand{\arraystretch}{1.15}\begin{tabular}{lr}
\shortstack[l]{Durchgang durch die Sonnennähe,\\Pariser Zeit 1807 September} .\;.\;.\; & $18{,}82718$\\
Neigung der Bahn .\;.\;.\;.\;.\;.\;.\;.\;.\;.\;.\;.\;.\;.\;.\;.\;.\; & $63^\circ\;14'\;28''1$\\
Aufsteigender Knoten .\;.\;.\;.\;.\;.\;.\;.\;.\;.\;.\;.\;.\;.\;.\;.\; & $266\;36\;51{,}7$\\
Sonnennähe .\;.\;.\;.\;.\;.\;.\;.\;.\;.\;.\;.\;.\;.\;.\;.\;.\;.\;.\;.\; & $271\;6\;7{,}5$\\
\multicolumn{2}{l}{(beide vom mittlern Aequinoctium gezählt)}\\
Logarithm des kleinsten Abstandes .\;.\;.\;.\;.\;.\;.\;.\;.\;.\;.\;.\; & $9{,}8122168$\\
Logarithm der mittlern Bewegung .\;.\;.\;.\;.\;.\;.\;.\;.\;.\;.\;.\; & $0{,}2418031$\\
\end{tabular}\end{center}


\par\indent Diese Elemente lassen bei den Breiten zu Anfang Decembers noch eine Abweichung von Einer
Minute zurück: durch Vertheilung würde sich dieser Fehler noch beträchtlich vermindern
lassen. Hr. \textsc{Bessel} hat indess lieber, mit Beisetzung der parabolischen Hypothese, die
Beobachtungen durch eine elliptische Bahn möglich genau vereinigen wollen, und wenn es
gleich in der Natur der Sache liegt, dass die Dimensionen derselben keiner scharfen
Bestimmung fähig sind, so scheint hieraus doch ziemlich zuverlässig hervorzugehen, dass
die Bahn wirklich elliptisch, und nicht hyperbolisch ist. Es gibt wenig Cometen, wo man
auch nur so viel mit entschiedener Gewissheit nach Einer Erscheinung behaupten kann. Die
Elemente der von Hrn. \textsc{Bessel} bestimmten elliptischen Bahn sind folgende:


\begin{center}\fontsize{9}{11}\selectfont\setlength{\tabcolsep}{3pt}\renewcommand{\arraystretch}{1.15}\begin{tabular}{lr}
Durchgangszeit durch die Sonnennähe September .\;.\;.\;.\;.\;.\; & $18{,}74986$\\
Neigung der Bahn .\;.\;.\;.\;.\;.\;.\;.\;.\;.\;.\;.\;.\;.\;.\; & $63^\circ\;10'\;53''2$\\
Aufsteigender Knoten .\;.\;.\;.\;.\;.\;.\;.\;.\;.\;.\;.\;.\;.\; & $266\;46\;3{,}1$\\
Sonnennähe .\;.\;.\;.\;.\;.\;.\;.\;.\;.\;.\;.\;.\;.\;.\;.\;.\;.\; & $271\;56\;0{,}1$\\
\multicolumn{2}{l}{(beide vom mittlern Aequinoctium gezählt.)}\\
Logarithm des kleinsten Abstandes .\;.\;.\;.\;.\;.\;.\;.\;.\;.\;.\; & $9{,}8105558$\\
Logarithm der mittlern Bewegung .\;.\;.\;.\;.\;.\;.\;.\;.\;.\;.\; & $0{,}2442946$\\
Excentricität .\;.\;.\;.\;.\;.\;.\;.\;.\;.\;.\;.\;.\;.\;.\;.\;.\;.\;.\; & $0{,}9958626$\\
Halbe grosse Axe .\;.\;.\;.\;.\;.\;.\;.\;.\;.\;.\;.\;.\;.\;.\;.\;.\;.\; & $156{,}253$\\
Umlaufszeit .\;.\;.\;.\;.\;.\;.\;.\;.\;.\;.\;.\;.\;.\;.\;.\;.\;.\;.\;.\; & $1953{,}2\text{ Jahre.}$\\
\end{tabular}\end{center}
